\src{180}
Les accouplements des deux colonnes de droite sont déduits de l'accouplement de
complexes de faisceaux
\[
\symup{\Homsh_{\Osh_X}(L_{*},\underline{\Omega}_X^{p})
\times\Homsh_{\Osh_X}(L'_{*},\underline{\Omega}_X^{p'})
\longrightarrow
\Homsh_{\Osh_X}(L_{*}\otimes L'_{*},\underline{\Omega}_X^{p+p'})}
\]
que l'on définit à l'aide du produit extérieur
$\symup{\underline{\Omega}_X^{p}\times\underline{\Omega}_X^{p'}\longrightarrow\underline{\Omega}_X^{p+p'}}$.
L'accouplement de la première colonne est le cup-produit (relatif au produit
extérieur). Je dis que la dernière ligne de (4.5) s'identifie au diagramme
d'homomorphismes analogue à (4.3), où $\symup{Y}$ est remplacé par $\symup{Y\cap Y'}$ et $\symup{p}$ par
$\symup{p+p'}$. Pour ceci, il suffit de montrer que $\symup{L\otimes L'}$ est une résolution
(évidemment localement libre) de $\symup{\Osh_{Y\cap Y'}}$. Or on a en effet
\[
  \symup{H_0(L\otimes L')=\Osh_Y\otimes\Osh_{Y'}=\Osh_{Y\cap Y'}}
\]
et
\[
  \symup{H_i(L\otimes L')=\operatorname{Tor}_{i}^{\Osh_X}(\Osh_Y,\Osh_{Y'})=0}
\]
pour $\symup{i>0}$, du fait que $\symup{Y}$ et $\symup{Y'}$ se coupent transversalement. Le théorème 1 résulte
maintenant de la formule :
\begin{equation}\tag{4.6}
  \symup{s_Y.s_{Y'}=s_{Y.Y'}}
\end{equation}
(où le produit du premier membre est celui de la dernière colonne de (4.5)).
Cette formule (4.6), de nature purement locale, se vérifie sans difficultés en
prenant pour $\symup{L_{*}}$ et $\symup{L'_{*}}$ les résolutions envisagées dans la proposition
4. On démontre de même (par une démonstration plus facile) que
$\symup{Z\longrightarrow P_X(Z)}$ est compatible avec le produit cartésien :
\begin{equation}\tag{4.7}
  \symup{P_{X\times X'}(Z\times Z')=P_X(Z)\otimes P_{X'}(Z')}
\end{equation}
(formule valable si $\symup{Z}$ resp. $\symup{Z'}$, est un cycle non singulier sur la variété
non singulière $\symup{X}$, resp. $\symup{X'}$, $\symup{Z\times Z'}$ étant considéré comme un cycle non
singulier sur $\symup{X\times X'}$). De (4.4) et (4.7) on déduit que $\symup{P_X(Z)}$ est aussi
compatible avec l'opération \fquote{image inverse} par un morphisme de variétés non
singulières $\symup{f:X\longrightarrow X'}$ :
\begin{equation}\tag{4.8}
  \symup{P_X\bigl(f^{-1}(Z')\bigr)=f^{*}\bigl(P_{X'}(Z')\bigr)}
\end{equation}
formule valable si $\symup{Z'}$ est un cycle non singulier sur $\symup{X'}$ tel que $\symup{f}$ soit
\fquote{transversale} à $\symup{Z'}$, i.e. tel que le graphe de $\symup{f}$ soit transversal au cycle
$\symup{X\times Z'}$ dans $\symup{X\times X'}$.

